\section{Connections}\label{sec:bridges}

This section records where the programme meets existing mathematics, the typed matchings proposed between the author's terms and exact objects, the $\mathbb{F}_1$ context, and the overlaps with the author's other workbenches. Each item says whether it is a theorem, an identity of objects, an analogy, or a proposal.

\subsection{Meetings with existing theory}

\begin{itemize}
\item \emph{Bost--Connes.} The normalised clock operators $S_n=\sqrt n\,V_n$ satisfy $S_mS_n=S_{mn}$ and $S_m^*S_n=S_{n/d}S_{m/d}^*$ ($d=\gcd(m,n)$), the relations of the semigroup crossed product $C(\widehat{\ZZ})\rtimes\NN^\times$ (Laca--Raeburn \cite{LacaRaeburn}), so the author's global quotient of all prime clocks is the space on which the Bost--Connes system acts (\lbl{P14}; \note{01\_} §4; theorem). The partition function of that system is $\zeta(\beta)$, with its phase transition at the pole $\beta=1$ \cite{BostConnes}; the link with the pole-at-$1$ ledger of Proposition~\ref{prop:poles} is an interpretation (no map is constructed).
\item \emph{Beurling.} The timed-prime reconstruction is Beurling's generalised-number construction $D=\exp_*(R)$ specialised to the actual primes (\lbl{TP9--TP13}; identity).
\item \emph{Arithmetic congruence monoids.} At $a=1/q$ the copy's Theorem~E is the factoriality criterion for $\{n\equiv1\bmod q\}$, a Krull monoid with class group $(\ZZ/q)^\times$, and the even-sheet test is the same criterion after dividing the class group by $\pm1$, which coincides with half-factoriality of the one-sided monoid, the Krull-monoid analogue of Carlitz's class-number-two theorem (\cite{BaginskiChapman,Carlitz}; Section~\ref{sec:sheets}; theorem).
\item \emph{Nyman--Beurling in three topologies.} In $\B$ the cover discrepancies generate the zero-jet ideal unconditionally, already for the degrees $\{2^a3^b\}$ when $t>0.0606$ (Proposition~\ref{prop:NB}); in $L^2(0,\infty)$ two-sided tests are dense unconditionally (Wiener's $L^2$ Tauberian theorem \cite{Wiener}); one-sided density in $L^2(0,1)$ is equivalent to RH (Beurling \cite{Beurling}, Báez-Duarte \cite{BaezDuarte}; Burnol's co-Poisson theory \cite{Burnol}), and Noor's Hardy-space form is also equivalent to RH \cite{Noor}. The three settings differ in topology, support and test class, and RH is the statement about the third (theorems; \note{09\_} §4(c), \note{16\_}).
\item \emph{Connes and Meyer.} Connes' §VIII harmonic distribution, done with the full Gram matrix, is the programme's \lbl{FGR}; Connes' weighted Hilbert quotients \cite{Connes} and Meyer's Fréchet quotient \cite{Meyer} are compared by one map whose kernel is the off-line jets (register S26, S28; theorems). Spectral synthesis on $Q$ is a local-description theorem in Krasichkov-Ternovskii's sense \cite{Krasichkov}; the positive transfer forms before the quotient are the Bochner--Schwartz theorem in Mellin coordinates \cite{Vladimirov}; the divisor potential belongs to the Littlewood-lemma family \cite{BalazardSaiasYor,SekatskiiBeltraminelliMerlini} (identities of objects and theorems, as stated in the programme and audited).
\item \emph{Interpolation.} The convergence of the primary decomposition (Theorem~\ref{thm:primary}) is the same kind of question as Berenstein--Taylor's interpolating varieties \cite{BerensteinTaylor} and Bondarenko--Radchenko--Seip's Fourier interpolation at the zeta zeros \cite{BondarenkoRadchenkoSeip}; the Gonek--Hejhal conjectures concern the averages $\sum|\zeta'(\rho)|^{-2k}$, and Theorem~\ref{thm:primary} asks the pointwise question (an analogy of questions; the theorem itself is proved).
\item \emph{Weil and Turing.} The mirror through the pole at $0$ followed by the passage between the lines is Weil's reflection through $\tfrac12$ ($T^{-1}A=\#$); on the diagonal of $T$-even data the chiral cross form is twice Weil's form (Proposition~\ref{prop:chiral}); the count of nontrivial holonomies plus $\sum\lfloor m_\rho/2\rfloor$ over the on-line zeros is Turing's index $(N(T)-V(T))/2$ (Proposition~\ref{prop:index}) (theorems).
\item \emph{Deligne.} Weil~II~§2.1 specialised to $\zeta$ is the classical proof of $\zeta(1+it)\ne0$, and the timed-prime measure maps onto Deligne's positive measure (\lbl{DB9}, \lbl{DR0--DR4}; theorems). In claude-ab's assessment the programme's even-power computations supply the Euler-product half of Deligne's §1.5 mechanism, and the missing half is the uniform pole bound (Section~\ref{sec:deligne}). In the programme it is ramification, not counting, that satisfies a relation of the form of Deligne's monodromy relation $FNF^{-1}=Q^{-1}N$ (as $P_bNP_b^{-1}=b^{-1}N$ on the image of $P_b$), and every solution of the relation on a jet block forces the weight shift $-2j$ of the programme's regrading; since such relations have solutions for every nilpotent, this is a comparison of relations only (\note{35\_} §3). Under that regrading the jet block of a zero of multiplicity $m$ satisfies the conclusion of Weil~II~1.8.4 with centre $2\re\rho-(m-1)$; for $m\ge2$ the block is mixed in Weil~II's sense (1.7.4), and ``pure'' only in the later terminology for Weil--Deligne representations \cite{TaylorYoshida}; the regrading is chosen, not derived, so this is bookkeeping (\note{36\_} Remark 36.R1).
\item \emph{Squares.} Deligne's step $\alpha\mapsto\alpha^2$ ($H^1\otimes H^1\hookrightarrow H^2$) has as additive shadow Kurokawa's tensor product of zeta functions, with spectral parameters $\rho+\rho'$ (Manin \cite{Manin}, §1.4); Weil's proof uses $\bar C\times\bar C$, and Connes proposes the square of the arithmetic site as its analogue \cite{ConnesSquare}. These are analogies, not identities (\note{08\_} §§3--4).
\item \emph{Connes--Consani.} The programme's doubled-source cohomology is Connes--Consani's §5 \cite{CCP1} pulled back (\lbl{DCP7.2}; \note{39\_}); the programme's adelic periodization complex is formal: after degree-zero coinvariants and in the algebraic representation category, $C\simeq\B_{\mathrm{pr}}[1]\oplus Q[0]$ equivariantly (\lbl{OMS7A.1}; \note{40\_} §4; theorems); the absolute twistor line \cite{CCtwistor} is linked to the CM elliptic curve $y^2=x^3-x$ through its $\mathbb{F}_{1^2}$-points (Section~\ref{sec:other}); the prime knots of \cite{CCknots} carry the conjugate character when reversed, and for primitive $\chi$ the functional equation puts the anti-character on the other side at the phase $W(\chi)$, with $W(\chi)W(\bar\chi)=1$ (\note{22\_} Proposition 22.8; theorems, with the reading as anti-numbers a proposal).
\item \emph{Two-prime density.} \lbl{MCL3.4}, \lbl{IH6} and \lbl{PL3} rest on one fact, that $(\log p)\ZZ+(\log q)\ZZ$ is dense in $\R$, the multiplicative form of the two-period phenomenon of mean-periodic functions \cite{SchwartzMP} (an observation about the programme's proofs; the density holds because $\log p/\log q$ is irrational, and the link with mean-periodic functions is a comparison of mechanisms).
\end{itemize}

\subsection{The anomaly dictionary (a proposal)}

The author's physical picture is matched term by term with exact objects (\note{22\_} §1). The objects and the facts proved about them hold independently of RH; the matching itself is a proposal.
\begin{center}\small
\begin{tabularx}{\linewidth}{@{}l X@{}}
\toprule
Author's term & Exact object\\
\midrule
anti-number & the clock run backwards with conjugation, $g\mapsto\cl{g(1/u)}/u$ (Mellin shadow $s\mapsto1-\bar s$)\\
chirality & the mirror $A(s)=-\bar s$ through the pole at $0$\\
the doublet & the orbit $\{\rho,\rho^\#,\rho^\#-1,\rho-1\}$ of an off-line zero in the two-line system\\
$\ZZ_2$ anomaly & the two lines as the two sheets of one double cover, deck transformation the central $-1$ of $SU(2)$\\
$\ZZ_1$ anomaly & the root-number phase $W(\chi)$, cancelled by the anti-sector: $W(\chi)W(\bar\chi)=1$\\
inflow & the bulk between the lines, which holds exactly the off-line zeros with $\re\rho<\tfrac12$ and their mirrors\\
the support $\tau$ & the pivot of the swing between $0$ and $1$: on the zeta side the reflection $s\mapsto1-\bar s$ about $\tfrac12$\\
\bottomrule
\end{tabularx}
\end{center}
In these terms RH has three exact forms: every zero is its own anti-number; the bulk is empty; and the Weil functional is nonnegative on every annihilation $g*g^\#$ (Weil's criterion; the annihilation's value at the unit, $\int|g|^2du$, is positive with or without RH).

\subsection{The \texorpdfstring{$\mathbb{F}_1$}{F1} context}

\begin{itemize}
\item The author's $\tau$ ``with $\ZZ_1$ and no $\ZZ_2$'' is realised at Connes--Consani's generic point $\eta$: every homeomorphism fixes $\eta$ and no exchanged label can sit there, while the invariant local data at $\eta$ still contain the sign $\varepsilon$. So ``no $\ZZ_2$'' has two exact readings: no exchanged label (which holds at $\eta$), or no sign at all ($\varepsilon=1$, which holds among field-valued points exactly in characteristic $2$, through the retraction $\mathbb{F}_{1^2}\to\mathbb{F}_1$).
\item The labels $\ZZ_n$ follow the principle that an $\mathbb{F}_{1^n}$-structure is a free $\ZZ/n$-action (Kapranov--Smirnov, via Connes--Consani--Marcolli \cite{CCMfun}, §3; a reading). The points of $\mathbb{P}^1$ over $\mathbb{F}_{1^2}$ are $\{0,\infty,\pm1\}$, and the branched double cover at them is $y^2=x^3-x$.
\item The Eulerian sheets $a=1/q$ of the even lattice family are exactly those with $(\ZZ/q)^\times=\{\pm1\}$, the units of $\mathbb{F}_{1^2}$, and exactly those on which the one-sided monoid is half-factorial (Proposition~\ref{prop:evensheets}); reading this as $\mathbb{F}_1\to\mathbb{F}_{1^2}$ is an interpretation.
\item The global quotient of all prime clocks is $\widehat{\ZZ}$ (Proposition~\ref{prop:clocks}), the Bost--Connes space, which Connes--Consani treat as an $\mathbb{F}_1$ object (a reading); the split-zero semiring $G(\ZZ)$ adds a Boolean point below the generic point of $\operatorname{Spec}\ZZ$, where the sign is forgotten, while Connes--Consani add a closed point $\infty$ glued along the generic point (a reading).
\item The question the weight lane arrived at, stated without the programme's vocabulary, is whether the reconstructed arithmetic carries (a) a square with a Künneth map $H^1\otimes H^1\to H^2$ and an auxiliary object with discrete weights, or (b) an independent cohomological source for a pole bound uniform in the tensor power (Section~\ref{sec:deligne}); the square is Manin's question of ``absolute Descartes powers'' of $\operatorname{Spec}\ZZ$ \cite{Manin}.
\end{itemize}

\subsection{Overlaps with the author's other workbenches}\label{ssec:overlaps}

The author maintains workbenches on the Erdős--Straus conjecture (ES), Yang--Mills (YM, which also holds the Navier--Stokes (NS), $S^6$ and Jacobian-map material), Erdős problem 817, and a Collatz workbench. The edges below are a selection of those located in these repositories (the Collatz side of the last row is an external record, the Clankers' Zenodo record 19900461; the author's Collatz workbench itself was not cloned). The column ``status'' says what this audit established. ``Audited'' means read and checked in \note{21\_}, \note{28\_} or \note{34\_}; ``located'' means found and described, not checked. Separate readers of the other workbenches are planned; this table is a selection of the zeta side of their intersection.

\begin{center}\small
\begin{longtable}{@{}p{0.2\linewidth}p{0.47\linewidth}p{0.27\linewidth}@{}}
\toprule
Edge & Shared object or lemma & Status\\
\midrule
\endhead
NS $\to$ ES & the torus endomorphism $J=\begin{psmallmatrix}3&1\\1&5\end{psmallmatrix}$ of the NS manuscript, used as an operator in 21 ES statements; exact aliasing on finite character groups; sharp constant $1/\sqrt{4+2\sqrt2}$ for the $\sqrt2$ directions & audited, correct (register S40); joint averages need the complete preimage (N-X)\\
NS $\to$ YM & NS velocity as a reducible $SU(2)$ connection; $-2\sum\operatorname{tr}\mathcal{F}_{ij}^2=\lambda^2|\operatorname{curl}u|^2$ & audited; correct conditional on imported NS bounds and on YM sections not read; zero-quotient sequence only as the coupling tends to $0$\\
zeta $\to$ YM & the Gram sandwich from a relative column error & audited, correct (register S41)\\
Jacobian map $\to$ NS kinematics $\to$ YM & $U=DF^{-1}(V\circ F)$, the escaping curve, the material tensor; YM Theorem~14.2 & audited: Lemma~\ref{lem:keller}; Theorem~14.2 correct conditional on fixed-box weak-coupling limits, quotients equal to the free two-gluon threshold (N-X)\\
$S^6\to$ YM & the imaginary period block $B$ of the $S^6$ period matrix enters only through $D=\det B$ & audited, correct (register S53)\\
$S^6\to$ ES & the integral monodromies of the $(3,4,\infty)$ period system give decorated affine-torsion covers of the ES divisor atlas & located (new edge; not audited)\\
ES $\leftrightarrow$ zeta & the split-zero semiring $G(R)=R\sqcup\{\tau\}$ with its unique Boolean character & audited (the zeta proof checked by hand)\\
zeta $\to$ ES & ES family-C results use zeta \lbl{NG20}--\lbl{NG21} as theorem dependencies & located (a dependency, not re-proved in ES)\\
ES $\leftrightarrow$ zeta & a four-dimensional Keller map $P$ with $\det DP=-2$, transported into the split-zero weighted conductor; the shared file is byte-identical in both repositories & $\det DP=-2$ audited; the rest located\\
NS $\to$ zeta & the released NS field as arithmetic input (Mellin transports, residue at $s=-1$) & located; depends on the imported NS existence theorem, whose independent verification the zeta side records as unfinished\\
817 $\to$ zeta & finite word sources of problem 817 mapped into the split-zero supported quotient & located (typed maps written; no zeta statement is entered)\\
Collatz $\leftrightarrow$ ES $\leftrightarrow$ zeta & the support-idempotent homogenisation $(x,h)\mapsto(ax+bh,h)$ & located (same formula on all three sides)\\
\bottomrule
\end{longtable}
\end{center}

One Collatz item is elementary and can be checked here: with $T(n)=4n+1$, $\theta(n)=3n+1$ and the odd step $U(n)=(3n+1)/2^{v_2(3n+1)}$, one has $\theta(T(n))=12n+4=4\theta(n)$, hence $U(T^jn)=U(n)$ for all $j$ (ES results bench, R05). The Erdős--Straus workbench states that it does not claim a proof or counterexample of the conjecture, and the Collatz workbench that it is not a claim that the Collatz conjecture has been solved; the Yang--Mills and Navier--Stokes material states its own scope limits, recorded in N-X.
